% GENERATED FILE. Edit canonical lessons, then run npm run generate:books.
% canonical-source-hash: fnv1a64:8a14e2f4b3ba8822
% canonical-lessons: RU-C82-this, RU-C82-this-2, RU-C82-that, RU-C82-that-2, RU-C82-such

\chapter{This, This, That, That, Such}
\label{ch:ru-this-this-that-that-such}
\hlchaptermodality{82}

\begin{chapteropening}
\textbf{By the end of this chapter:} \emph{I can say this, this, that, that and such.}

Complete the last lesson of chapter 82: i can say this, this, that, that and such.
\end{chapteropening}

\section[étot]{\ru{этот} (étot) — this (masculine)}

\label{lesson:RU-C82-this}



\begin{quote}
Before the new one: say the Russian for a park, then the Russian for a taxi.
\end{quote}

\subsection*{You'll want to know: \ru{этот}}
\textbf{\ru{этот}} (\emph{étot}) — \textquotedblleft{}this (masculine)\textquotedblright{}.

\begin{grammarlens}[title={What you've built}]
One more word for this chapter.
\end{grammarlens}

\subsection*{Your turn}
\begin{itemize}
\raggedright
  \item \emph{Say it:} \emph{étot}
  \item \emph{Say it:} \emph{étot}, once more
  \item \emph{Say it:} say \emph{étot}
  \item \emph{Recall:} say the Russian for a park, then the Russian for a taxi, then say \emph{étot} again
\end{itemize}

\begin{tcolorbox}[breakable,colback=teal!4,colframe=teal!35!black,title={Before you move on}]
What does \ru{этот} mean? (\textquotedblleft{}This (masculine)\textquotedblright{}.) Say it once more.
\end{tcolorbox}
\section[éta]{\ru{эта} (éta) — this (feminine)}

\label{lesson:RU-C82-this-2}



\begin{quote}
Before the new one: say the Russian for a taxi, then the Russian for this.
\end{quote}

\subsection*{You'll want to know: \ru{эта}}
\textbf{\ru{эта}} (\emph{éta}) — \textquotedblleft{}this (feminine)\textquotedblright{}.

\begin{grammarlens}[title={What you've built}]
One more word for this chapter.
\end{grammarlens}

\subsection*{Your turn}
\begin{itemize}
\raggedright
  \item \emph{Say it:} \emph{éta}
  \item \emph{Say it:} \emph{éta}, once more
  \item \emph{Say it:} say \emph{éta}
  \item \emph{Recall:} say the Russian for a taxi, then the Russian for this, then say \emph{éta} again
\end{itemize}

\begin{tcolorbox}[breakable,colback=teal!4,colframe=teal!35!black,title={Before you move on}]
What does \ru{эта} mean? (\textquotedblleft{}This (feminine)\textquotedblright{}.) Say it once more.
\end{tcolorbox}
\section[tot]{\ru{тот} (tot) — that (masculine)}

\label{lesson:RU-C82-that}



\begin{quote}
Before the new one: say the Russian for this, then the Russian for this.
\end{quote}

\subsection*{You'll want to know: \ru{тот}}
\textbf{\ru{тот}} (\emph{tot}) — \textquotedblleft{}that (masculine)\textquotedblright{}.

\begin{grammarlens}[title={What you've built}]
One more word for this chapter.
\end{grammarlens}

\subsection*{Your turn}
\begin{itemize}
\raggedright
  \item \emph{Say it:} \emph{tot}
  \item \emph{Say it:} \emph{tot}, once more
  \item \emph{Say it:} say \emph{tot}
  \item \emph{Recall:} say the Russian for this, then the Russian for this, then say \emph{tot} again
\end{itemize}

\begin{tcolorbox}[breakable,colback=teal!4,colframe=teal!35!black,title={Before you move on}]
What does \ru{тот} mean? (\textquotedblleft{}That (masculine)\textquotedblright{}.) Say it once more.
\end{tcolorbox}
\section[ta]{\ru{та} (ta) — that (feminine)}

\label{lesson:RU-C82-that-2}



\begin{quote}
Before the new one: say the Russian for this, then the Russian for that.
\end{quote}

\subsection*{You'll want to know: \ru{та}}
\textbf{\ru{та}} (\emph{ta}) — \textquotedblleft{}that (feminine)\textquotedblright{}.

\begin{grammarlens}[title={What you've built}]
One more word for this chapter.
\end{grammarlens}

\subsection*{Your turn}
\begin{itemize}
\raggedright
  \item \emph{Say it:} \emph{ta}
  \item \emph{Say it:} \emph{ta}, once more
  \item \emph{Say it:} say \emph{ta}
  \item \emph{Recall:} say the Russian for this, then the Russian for that, then say \emph{ta} again
\end{itemize}

\begin{tcolorbox}[breakable,colback=teal!4,colframe=teal!35!black,title={Before you move on}]
What does \ru{та} mean? (\textquotedblleft{}That (feminine)\textquotedblright{}.) Say it once more.
\end{tcolorbox}
\section[takóy]{\ru{такой} (takóy) — such, like this}

\label{lesson:RU-C82-such}



\begin{quote}
Before the new one: say the Russian for that, then the Russian for that.
\end{quote}

\subsection*{You'll want to know: \ru{такой}}
\textbf{\ru{такой}} (\emph{takóy}) — \textquotedblleft{}such, like this\textquotedblright{}.

\begin{grammarlens}[title={What you've built}]
One more word for this chapter.
\end{grammarlens}

\subsection*{Your turn}
\begin{itemize}
\raggedright
  \item \emph{Say it:} \emph{takóy}
  \item \emph{Say it:} \emph{takóy}, once more
  \item \emph{Say it:} say \emph{takóy}
  \item \emph{Recall:} say the Russian for that, then the Russian for that, then say \emph{takóy} again
\end{itemize}

\begin{tcolorbox}[breakable,colback=teal!4,colframe=teal!35!black,title={Before you move on}]
What does \ru{такой} mean? (\textquotedblleft{}Such, like this\textquotedblright{}.) Say it once more.
\end{tcolorbox}
